\chapter{Introduction}\label{ch:introduction}

This is the template to use for your master thesis Software Engineering.
It provides you with a structure that we expect you to follow, with guidelines on the contents of each section.
Depending on your research you can deviate from this structure, usually by renaming or adding some chapters in the middle.
Do discuss this with your academic supervisor first!
Be sure to change the variables under the \verb|Options| banner set in \verb|main.tex|, such as title, subtitle, author, supervisor, company, \etc{}

The Introduction chapter itself is to provide your reader with the context of your research, the problem you are trying to solve.

\section{Context}
Provide the context of your research, and explain the importance and relevance of your topic within the field of software engineering.

\section{Problem Statement}
Give a concise description of the challenge you aim to address.
The goal is to define the problem in clear terms, establish its significance, and provides a foundation for developing objectives, methods, and potential solutions.
Be specific and avoid vague language. 

\section{Research Questions}
Explicitly list the research questions or hypotheses you address in your research.
The research questions should directly address the problem statement, as defined in the previous section, and guide your research methodology, as described in the next section.
It is helpful to enumerate them, for example as \emph{RQ1}, \emph{RQ2}, so you can refer to them in later chapters. 

\section{Approach}
Give a brief overview of the approach you will take to answer the research questions.
Define the research methods you used (\eg{} case study, systematic literature review, action research, \etc{}), and the general steps you will take to conduct your research.
This will provide the reader with a high-level roadmap.

\section{Outline}
Give the reader a short overview of the structure of your thesis, and what to find in each chapter.